\begin{proofbox}
	\textbf{\textcolor{CProof}{思路提示}}

	通过通分 $ \frac1a+\frac1b = \frac{a+b}{ab} $，将倒数和化为积的倒数；再利用基本不等式由 $a+b=1$ 得到 $ab\le\frac14$，从而 $\frac1{ab}\ge4$。

	\medskip
	\textbf{\textcolor{CProof}{正式推导}}
	\begin{description}[style=nextline, leftmargin=2.8em, labelsep=0.5em, itemsep=0.55em]
		\item[\textbf{步骤一（通分合并）：}]
			将两个倒数通分：
			\[
				\begin{aligned}
					\frac1a+\frac1b &= \frac{b}{ab} + \frac{a}{ab} \\
					&= \frac{a+b}{ab}.
			\end{aligned}
			\]
			\par\textbf{理由：} 分母不相同，通分后合并便于利用条件。
		\item[\textbf{步骤二（代入条件）：}]
			代入条件 $a+b=1$：
			\[
				\frac1a+\frac1b = \frac{1}{ab}.
			\]
			\par\textbf{理由：} 简化表达式，将倒数和转化为积的倒数。
		\item[\textbf{步骤三（基本不等式取积范围）：}]
			由 $a>0,\ b>0$ 及 $a+b=1$，使用基本不等式：
			\[
				\begin{aligned}
					ab &\le \left(\frac{a+b}{2}\right)^2 \\
					&= \left(\frac{1}{2}\right)^2 \\
					&= \frac14.
			\end{aligned}
			\]
			\par\textbf{理由：} 和为定值时，乘积在两者相等时取最大值。
		\item[\textbf{步骤四（推导倒数不等式）：}]
			由 $ab\le\frac14$ 且 $ab>0$，取倒数得：
			\[
				\frac1{ab} \ge 4.
			\]
			\par\textbf{理由：} 倒数函数在正数区间上是反比例递减，因此积越大倒数越小，故不等式方向反转。
		\item[\textbf{步骤五（取等条件）：}]
			基本不等式取等条件为 $a=b$，结合 $a+b=1$ 得 $a=b=\frac12$。此时 $\frac1a+\frac1b=4$。
			\par\textbf{理由：} 等号必须同时满足基本不等式取等和条件等式。
	\end{description}

	\medskip
	\textbf{\textcolor{CProof}{结论回扣}}

	综合以上步骤，当 $a>0,\ b>0$ 且 $a+b=1$ 时，有 $\frac1a+\frac1b\ge4$，等号当且仅当 $a=b=\frac12$ 成立。原结论成立。
\end{proofbox}
